% Vietnamese translation for OI Public Library
% Source-File: IOI中国国家候选队论文/2004/许智磊.ppt
\documentclass[11pt]{article}
\usepackage{oipl-vn}

\newcommand{\Oh}{\mathcal{O}}
\newcommand{\SA}{\operatorname{SA}}
\newcommand{\Rank}{\operatorname{Rank}}
\newcommand{\LCP}{\operatorname{LCP}}

\title{Mảng hậu tố\\{\large Ghi chú theo slide}}
\author{Xu Zhilei}
\date{IOI 2004}

\begin{document}
\maketitle

\origtitle{Suffix array}
\sourcepath{IOI 2004 / Xu Zhilei.ppt}

\section{Khái niệm}

Các slide đi từ những ký hiệu cơ bản: độ dài \(\operatorname{len}(S)=n\), ký
tự \(S[i]\), xâu con \(S[i..j]\), hậu tố \(S[i..n]\) và
\(\operatorname{Suffix}(i)\). Mảng hậu tố \(\SA\) là thứ tự tăng dần từ điển
của mọi hậu tố; mảng \(\Rank[i]\) là vị trí của hậu tố bắt đầu tại \(i\) trong
thứ tự đó.

\section{Xây dựng}

Cách thô so sánh trực tiếp mọi hậu tố có thể tới \(\Oh(n^2)\) hoặc tệ hơn
tùy cách sắp xếp. Bài trình bày giới thiệu thuật toán nhân đôi: ở vòng
\(k\), mỗi hậu tố được đặc trưng bởi cặp hạng của hai khối độ dài \(2^k\).
Sau khi sắp xếp các cặp hạng, ta thu được hạng cho khối độ dài \(2^{k+1}\).
Lặp đến khi mọi hạng phân biệt cho độ phức tạp \(\Oh(n\log n)\) nếu dùng sắp
xếp tuyến tính theo hạng.

\section{Mảng LCP và ứng dụng}

Mảng \(\LCP\) ghi độ dài tiền tố chung dài nhất của hai hậu tố kề nhau trong
\(\SA\). Kết hợp \(\SA\), \(\Rank\) và \(\LCP\), ta có thể so khớp mẫu, tìm
xâu con lặp, hoặc xử lý bài xâu đối xứng bằng cách chuyển truy vấn LCP thành
truy vấn cực tiểu trên đoạn.

\section{Thông điệp}

Mảng hậu tố là lựa chọn thực dụng trong thi đấu vì dùng mảng thuần, ít con
trỏ hơn cây hậu tố, nhưng vẫn giữ được hầu hết sức mạnh cần thiết cho bài
toán xâu.

\end{document}
